    \chapter{Estudo Prévio}
\section{Análise das APIs de Inferência}

O estudo inicial comparou 3 APIs de inferência de utilização comum, sendo estas a \textbf{OpenAI} \cite{openai_api}, \textbf{Anthropic} \cite{anthropic_api} e \textbf{Gemini} \cite{google_gemini_api}, para este estudo, utilizou-se o \textbf{Ollama} para analisar a API da \textbf{Anthropic}, o \textbf{Groq} para a da \textbf{OpenAI}, e o \textbf{AI Studio} da própria Google para avaliar a API do \textbf{Gemini}. 

Este estudo não teve apenas um papel de descrever as APIs, mas sim servir como base e ponto de partida para observar os conceitos que podiam ser normalizados no \textbf{OmniAI SDK} e os pontos onde a normalização teria de preservar extensões específicas de cada fornecedor. Para esse efeito, foram executados e documentados pedidos focados apenas em formato textual de pedidos de geração simples, pedidos com \textit{streaming}, pedidos com \textit{tool calling}, respostas com contabilização de \textit{tokens} e casos de configuração de geração.

O estudo das APIs de inferência mostrou que estas partilham diversos conceitos semelhantes, mas permitem a utilização dos mesmos de maneira diferente. A \textbf{OpenAI} organiza a interação em torno de \texttt{chat/completions}, com \texttt{messages}, \texttt{choices}, \texttt{tool\_calls} e eventos incrementais baseados em \texttt{delta} \cite{openai_api}. A \textbf{Anthropic} usa uma API de \texttt{messages}, onde o \texttt{system prompt} está no nível de topo, o campo \texttt{max\_tokens} é obrigatório, e o conteúdo da resposta é uma lista de blocos como \texttt{text}, \texttt{tool\_use} ou \texttt{thinking} \cite{anthropic_api}. A \textbf{Gemini} modela o pedido como \texttt{contents} compostos por \texttt{parts}, agrupa parâmetros em \texttt{generationConfig}, representa ferramentas por \texttt{functionDeclarations} e introduz conceitos próprios como \texttt{thought signatures} e \texttt{thinkingConfig} \cite{google_gemini_api}. Ao longo deste capítulo as diferenças e semelhanças iram ser apresentadas em mais detalhe.

\subsection{Construção do Pedido}

\begin{table}[H]
    \centering
    \footnotesize
    \renewcommand{\arraystretch}{1.4}
    \begin{tabularx}{\textwidth}{|>{\raggedright\arraybackslash}p{2.8cm}|Y|Y|Y|}
        \hline
        \textbf{Construção do Pedido} & \textbf{OpenAI} & \textbf{Anthropic} & \textbf{Gemini} \\
        \hline
        \textbf{Autenticação} & Cabeçalho \texttt{Authorization: Bearer <TOKEN>}. & Cabeçalho \texttt{x-api-key} e obrigatoriedade de \texttt{anthropic-}\allowbreak\texttt{version}. & Cabeçalho \texttt{x-goog-api-key} ou parâmetro \texttt{?key=} no URI. \\
        \hline
        \textbf{URI (\textit{Endpoint})} & Fixo por tarefa (e.g., \texttt{.../v1/chat/}\allowbreak\texttt{completions}). & Fixo para mensagens (\texttt{.../v1/messages}). & Dinâmico, acoplando modelo e método (\texttt{.../models/}\allowbreak\texttt{\{model\}:\{method\}}). \\
        \hline
        \textbf{Identificação do Modelo} & Definido via propriedade \texttt{"model"} na raiz do \textit{payload} JSON. & Definido via propriedade \texttt{"model"} na raiz do \textit{payload} JSON. & Injetado dinamicamente no próprio URI de chamada \textbf{HTTP}. \\
        \hline
        \textbf{Estrutura do Histórico} & \textit{Array} \texttt{messages} contendo objetos com \texttt{role} e \texttt{content}. & \textit{Array} \texttt{messages} contendo objetos com \texttt{role} e \texttt{content}. & \textit{Array} \texttt{contents} contendo objetos com \texttt{role} e \textit{array} \texttt{parts}. \\
        \hline
        \textbf{Conteúdo da Mensagem} & O \texttt{content} pode ser \textit{string} simples ou \textit{array} polimórfico. & O \texttt{content} exige estritamente um \textit{array} de blocos estruturados. & O conteúdo divide-se num \textit{array} \texttt{parts} de objetos tipados. \\
        \hline
        \textbf{Instruções de Sistema} & Tratadas como mensagem inicial com \texttt{role: "system"}. & Declaradas na propriedade isolada \texttt{"system"} na raiz do \textit{payload}. & Declaradas no objeto isolado \texttt{"system\_}\allowbreak\texttt{instruction"} na raiz. \\
        \hline
        \textbf{Parâmetros de Geração} & Propriedades soltas na raiz (e.g., \texttt{temperature}, \texttt{max\_tokens}). & Propriedades soltas na raiz (\texttt{"max\_tokens"} é obrigatório). & Agrupados num único objeto configuracional (\texttt{"generation}\allowbreak\texttt{Config"}). \\
        \hline
        \textbf{Declaração de Ferramentas} & Propriedade \texttt{tools[].function} contendo \texttt{parameters}. & Propriedade \texttt{tools[]} contendo \texttt{input\_schema}. & Propriedade \texttt{tools[].}\allowbreak\texttt{functionDeclarations} contendo \texttt{parameters}. \\
        \hline
        \textbf{Ativação de \textit{Streaming}} & Propriedade \texttt{"stream": true} na raiz do corpo JSON. & Propriedade \texttt{"stream": true} na raiz do corpo JSON. & URI distinto (\texttt{streamGenerate}\allowbreak\texttt{Content}) com parâmetro \texttt{?alt=sse}; não existe campo no corpo. \\
        \hline
        \textbf{Configuração de Raciocínio} & Suporte nativo e automático na família de modelos \texttt{o1}/\texttt{o3}; sem campo dedicado de ativação. & Objeto \texttt{"thinking"} na raiz, com \texttt{type} e \texttt{budget\_tokens} para controlar o orçamento. & Objeto \texttt{thinkingConfig} dentro de \texttt{generationConfig}, com opções como \texttt{thinkingLevel} e \texttt{include\_thoughts}. \\
        \hline
        \textbf{Saídas Estruturadas} & Controladas via \texttt{"response\_}\allowbreak\texttt{format"} na raiz (e.g., \texttt{\{"type": "json\_object"\}}). & Implementadas forçando uma ferramenta dedicada via \texttt{tool\_choice} para extrair dados. & Definidas via \texttt{responseMimeType} e \texttt{responseJsonSchema} dentro de \texttt{generationConfig}. \\
        \hline
    \end{tabularx}
    \caption{Mapeamento da estrutura de pedidos e configuração (\textit{Requests}).}
    \label{tab:api-request-comparison}
\end{table}

    Com base nos pedidos executados durante o estudo, foi possível identificar os campos JSON mais relevantes de cada fornecedor. Na Figura~\ref{fig:json-request-skeletons} indentificou se as estruturas dessas mensagens, evidenciando as diferenças de organização e de nomenclatura entre os três fornecedores. Os exemplos JSON completos dos pedidos encontram-se no Apêndice~\ref{apx:json-requests}.

\begin{figure}[H]
    \centering
    \begin{minipage}[t]{0.32\textwidth}
        \centering
        \textbf{OpenAI Request}
\begin{lstlisting}[language=json, basicstyle=\ttfamily\scriptsize]
{
  "model": "gpt-4o",
  "messages": [
    {
      "role": "system",
      "content": "..."
    },
    {
      "role": "user",
      "content": "..."
    }
  ],
  "temperature": 0.7,
  "max_tokens": 150,
  "tools": [
    {
      "type": "function",
      "function": {
        "name": "func",
        "parameters": {...}
      }
    }
  ],
  "tool_choice": "auto",
  "stream": false
}
\end{lstlisting}
    \end{minipage}
    \hfill
    \begin{minipage}[t]{0.32\textwidth}
        \centering
        \textbf{Anthropic Request}
\begin{lstlisting}[language=json, basicstyle=\ttfamily\scriptsize]
{
  "model": "claude-3-5-sonnet",
  "system": "...",
  "messages": [
    {
      "role": "user",
      "content": [
        {
          "type": "text",
          "text": "..."
        }
      ]
    }
  ],
  "max_tokens": 150,
  "temperature": 0.7,
  "tools": [
    {
      "name": "func",
      "input_schema": {...}
    }
  ],
  "tool_choice": {
    "type": "auto"
  }
}
\end{lstlisting}
    \end{minipage}
    \hfill
    \begin{minipage}[t]{0.32\textwidth}
        \centering
        \textbf{Gemini Request}
\begin{lstlisting}[language=json, basicstyle=\ttfamily\scriptsize]
    URI: models/{model}:{method}
{
  
  "contents": [
    {
      "role": "user",
      "parts": [
        {"text": "..."}
      ]
    }
  ],
  "systemInstruction": {
    "parts": [
      {"text": "..."}
    ]
  },
  "generationConfig": {
    "temperature": 0.7,
    "maxOutputTokens": 150
  },
  "tools": [
    {
      "functionDeclarations": [
        {
          "name": "func",
          "parameters": {...}
        }
      ]
    }
  ]
}
\end{lstlisting}
    \end{minipage}
    \caption{Estruturas JSON esquemáticas e simplificadas de pedidos para cada fornecedor.}
    \label{fig:json-request-skeletons}
\end{figure}

A análise dos campos de pedido revela que, embora os três fornecedores partilhem conceitos equivalentes, cada um organiza o corpo JSON de forma distinta. A \textbf{OpenAI} \cite{openai_api} e a \textbf{Anthropic} \cite{anthropic_api} colocam a maioria dos parâmetros de configuração como propriedades soltas na raiz do objeto, enquanto a \textbf{Gemini} agrupa-os dentro de \texttt{generationConfig} \cite{google_gemini_api}. A diferença mais notável está na identificação do modelo. Neste caso, a \textbf{OpenAI} e a \textbf{Anthropic} incluem-na no corpo do pedido, enquanto a \textbf{Gemini} a incorpora diretamente no URI da chamada \textbf{HTTP}, tornando o \textit{endpoint} dinâmico. O \textit{system prompt} também diverge consideravelmente: a \textbf{OpenAI} trata-o como uma mensagem com o papel \texttt{system} \cite{openai_api}, a \textbf{Anthropic} eleva-o a uma propriedade de topo \cite{anthropic_api}, e a \textbf{Gemini} define-o através de um objeto \texttt{system\_instruction} separado \cite{google_gemini_api}. Estas diferenças tornam evidente que uma tradução superficial dos nomes dos campos seria insuficiente para construir uma abstração e um domínio comum que seja capaz de cobrir os três fornecedores de inferência.

\subsection{Processamento da Resposta}

\begin{table}[H]
    \centering
    \footnotesize
    \renewcommand{\arraystretch}{1.4}
    \begin{tabularx}{\textwidth}{|>{\raggedright\arraybackslash}p{2.8cm}|Y|Y|Y|}
        \hline
        \textbf{Processamento da Resposta} & \textbf{OpenAI} & \textbf{Anthropic} & \textbf{Gemini} \\
        \hline
        \textbf{Estrutura Raiz da Resposta} & Objeto \texttt{choices[]} contendo \texttt{message} com \texttt{role} e \texttt{content}. & Diretamente na raiz: \texttt{role}, \texttt{content[]}, \texttt{stop\_reason} e \texttt{usage}. & Objeto \texttt{candidates[]} contendo \texttt{content} com \texttt{role} e \texttt{parts[]}. \\
        \hline
        \textbf{Invocação pelo Modelo (\textit{Tool Call})} & Retorna \texttt{tool\_calls[]} com \texttt{arguments} em \textit{string} JSON serializada. & Retorna um bloco \texttt{type: "tool\_use"} com \texttt{input} em JSON nativo. & Retorna uma \texttt{part} contendo \texttt{functionCall} com \texttt{args} (JSON nativo). \\
        \hline
        \textbf{Resposta da Ferramenta} & Adicionada como nova mensagem com \texttt{role: "tool"} e \texttt{tool\_call\_id}. & Adicionada sob \texttt{role: "user"} num bloco \texttt{type: "tool\_result"} com \texttt{tool\_use\_id}. & Adicionada como nova mensagem contendo \texttt{functionResponse} com \texttt{name} e \texttt{response}. \\
        \hline
        \textbf{Razão de Paragem} & Campo \texttt{finish\_reason} com valores como \texttt{"stop"}, \texttt{"tool\_calls"} ou \texttt{"length"}. & Campo \texttt{stop\_reason} com valores como \texttt{"end\_turn"}, \texttt{"tool\_use"} ou \texttt{"max\_tokens"}. & Campo \texttt{finishReason} com valores como \texttt{"STOP"} ou \texttt{"MAX\_TOKENS"} (notação em maiúsculas). \\
        \hline
        \textbf{Saídas Estruturadas} & Controladas via propriedade \texttt{"response\_format"} na raiz. & Implementadas forçando uma ferramenta dedicada via \texttt{tool\_choice}. & Definidas via \texttt{response}\allowbreak\texttt{MimeType} e \texttt{response}\allowbreak\texttt{JsonSchema}. \\
        \hline
        \textbf{Comportamento \textit{Streaming}} & Ativado via \texttt{stream: true}; eventos \textbf{SSE} com formato \texttt{delta}. & Ativado via \texttt{stream: true}; eventos \textbf{SSE} fortemente tipados (\texttt{message\_start}, \texttt{content\_block\_delta}, etc.). & Consumido através de um \textit{endpoint} distinto com parâmetro \texttt{?alt=sse}; cada evento contém um objeto de resposta parcial completo. \\
        \hline
        \textbf{Metadados de Raciocínio} & Suporte nativo e automático na família de modelos \texttt{o1}/\texttt{o3}. & Ativado via \texttt{"thinking"}; devolvido num bloco \texttt{"thinking"} com \texttt{signature}. & Ativado via \texttt{thinking}\allowbreak\texttt{Config}; suporta entrega de \texttt{thought}\allowbreak\texttt{Signature} nos \texttt{parts}. \\
        \hline
        \textbf{Contabilização de \textit{Tokens}} & Objeto \texttt{usage} (campos \texttt{prompt\_tokens}, \texttt{completion\_}\allowbreak\texttt{tokens}, \texttt{total\_tokens}). & Objeto \texttt{usage} (campos \texttt{input\_tokens}, \texttt{output\_tokens}; sem total explícito). & Objeto \texttt{usageMetadata} (\texttt{prompt}\allowbreak\texttt{TokenCount}, \texttt{candidates}\allowbreak\texttt{TokenCount}, \texttt{totalTokenCount}, \texttt{thoughtsTokenCount}). \\
        \hline
        \textbf{Identificador da Resposta} & Campo \texttt{id} (e.g., \texttt{"chatcmpl-31"}) e metadados \texttt{object}, \texttt{created}. & Campo \texttt{id} (e.g., \texttt{"msg\_01Xf..."}) e campo \texttt{type: "message"}. & Campo \texttt{responseId} e \texttt{modelVersion} ao nível raiz da resposta. \\
        \hline
    \end{tabularx}
    \caption{Mapeamento técnico do processamento de respostas, execução e telemetria (\textit{Responses}).}
    \label{tab:api-response-comparison}
\end{table}

    À semelhança dos pedidos, a estrutura JSON das respostas foi esquematizada para evidenciar as difereças e semelhanças de cada resposta. A Figura~\ref{fig:json-response-skeletons} apresenta essa comparação estrutural. Os exemplos JSON completos das respostas encontram-se no Apêndice~\ref{apx:json-responses}.

\begin{figure}[H]
    \centering
    \begin{minipage}[t]{0.32\textwidth}
        \centering
        \textbf{OpenAI Response}
\begin{lstlisting}[language=json, basicstyle=\ttfamily\scriptsize]
{
  "id": "chatcmpl-...",
  "choices": [
    {
      "index": 0,
      "message": {
        "role": "assistant",
        "content": "...",
        "tool_calls": [
          {
            "id": "call_...",
            "type": "function",
            "function": {
              "name": "func",
              "arguments": "{\"...\"}"
            }
          }
        ]
      },
      "finish_reason": "stop"
    }
  ],
  "usage": {
    "prompt_tokens": 10,
    "completion_tokens": 20
  }
}
\end{lstlisting}
    \end{minipage}
    \hfill
    \begin{minipage}[t]{0.32\textwidth}
        \centering
        \textbf{Anthropic Response}
\begin{lstlisting}[language=json, basicstyle=\ttfamily\scriptsize]
{
  "id": "msg-...",
  "role": "assistant",
  "content": [
    {
      "type": "text",
      "text": "..."
    },
    {
      "type": "tool_use",
      "id": "toolu_...",
      "name": "func",
      "input": {...}
    }
  ],
  "stop_reason": "end_turn",
  "usage": {
    "input_tokens": 10,
    "output_tokens": 20
  }
}
\end{lstlisting}
    \end{minipage}
    \hfill
    \begin{minipage}[t]{0.32\textwidth}
        \centering
        \textbf{Gemini Response}
\begin{lstlisting}[language=json, basicstyle=\ttfamily\scriptsize]
{
  "candidates": [
    {
      "content": {
        "role": "model",
        "parts": [
          {"text": "..."},
          {
            "functionCall": {
              "name": "func",
              "args": {...}
            }
          }
        ]
      },
      "finishReason": "STOP"
    }
  ],
  "usageMetadata": {
    "promptTokenCount": 10,
    "candidatesTokenCount": 20
  }
}
\end{lstlisting}
    \end{minipage}
    \caption{Estruturas JSON esquemáticas e simplificadas de respostas para cada fornecedor.}
    \label{fig:json-response-skeletons}
\end{figure}

A comparação das respostas reforça a perceção de que as três APIs divergem significativamente na organização estrutural, apesar de exprimirem os mesmos conceitos. A interface da \textbf{OpenAI} encapsula a resposta dentro de um \textit{array} denominado \texttt{choices}, que contém um objeto \texttt{message} \cite{openai_api}; a interface da \textbf{Anthropic} expõe o conteúdo diretamente na raiz através de um \textit{array} \texttt{content} com blocos tipados \cite{anthropic_api}; e a interface da \textbf{Gemini} utiliza \texttt{candidates} com \texttt{parts} e conceitos adicionais, como \texttt{thoughtSignature} \cite{google_gemini_api}. 

Um aspeto particularmente relevante é a representação dos argumentos em chamadas de ferramenta: a \textbf{OpenAI} serializa-os como uma \textit{string} JSON dentro do campo \texttt{arguments}, obrigando a uma desserialização adicional \cite{openai_api}, enquanto a \textbf{Anthropic} e a \textbf{Gemini} devolvem JSON nativo nos campos \texttt{input} \cite{anthropic_api} e \texttt{args} \cite{google_gemini_api}, respetivamente. A contabilização de \textit{tokens} também difere na nomenclatura e granularidade: a API da \textbf{Gemini}, por exemplo, expõe um campo adicional \texttt{thoughtsTokenCount} que permite distinguir os \textit{tokens} de raciocínio dos \textit{tokens} de conteúdo efetivo \cite{google_gemini_api}.

\subsection{Divergências na Padronização de Esquemas de Validação}
\label{subsec:divergencias-esquemas}

Para além das diferenças do corpo JSON, é importante referir a divergência na forma como cada fornecedor padroniza a tipagem de dados estruturados, nomeadamente na declaração de ferramentas (\textit{tools}) e nas saídas estruturadas (\textit{structured outputs}). 

Enquanto a \textbf{OpenAI} \cite{openai_structured_outputs} e a \textbf{Anthropic} \cite{anthropic_tool_use} se apoiam nativamente na norma \textit{JSON Schema} \cite{json_schema_spec} para definir propriedades e parâmetros, a API da \textbf{Gemini} adota um subconjunto da especificação \textit{OpenAPI 3.0} \cite{openapi_3_spec} (através do seu objeto \texttt{Schema} \cite{gemini_schema_object}). Isto reforça que a construção de uma camada de abstração exige não só o mapeamento de chaves, mas também a tradução e compatibilização entre diferentes definições de esquemas.

\subsection{Conclusão do Estudo das APIs de Inferência}

  A realização deste estudo envolveu dificuldades práticas que merecem ser documentadas. Para a API da \textbf{OpenAI}, o acesso direto à API oficial implica custos de utilização, pelo que os testes foram realizados utilizando a API da \textbf{Groq} \cite{groq_api}, que oferece um plano gratuito e implementa uma interface compatível com o formato da \textbf{OpenAI} (\texttt{/v1/chat/completions}). Adicionalmente, foram executados testes complementares contra uma instância local de \textbf{Ollama} \cite{ollama_docs}, que também expõe um \textit{endpoint} compatível com o formato \textbf{OpenAI}. Para a API da \textbf{Anthropic}, a situação foi semelhante: a API oficial requer um plano de subscrição pago, pelo que os testes foram realizados utilizando o \textbf{Ollama}, que disponibiliza um \textit{endpoint} compatível com o formato da \textbf{Anthropic} (\texttt{/v1/messages}). Embora o \textbf{Ollama} implemente corretamente as interfaces de compatibilidade para pedidos simples, de \textit{streaming} e com parâmetros de configuração, algumas funcionalidades mais avançadas, nomeadamente o \textit{tool calling} estruturado e a geração de blocos de raciocínio (\textit{thinking}) apresentaram limitações quando executadas em modelos de pequena dimensão, como o \texttt{llama3.2:1b}. Nesses casos, os exemplos de resposta foram ajustados para refletir com fidelidade o comportamento documentado das APIs oficiais. A API do \textbf{Gemini} foi a única cujos testes foram executados diretamente contra a API oficial, nomeadamente através do \textbf{AI Studio} \cite{google_aistudio}, uma vez que este disponibiliza um nível de acesso gratuito.

  Do ponto de vista técnico, o estudo permitiu concluir que as três APIs partilham um conjunto de conceitos fundamentais, como as mensagens, papéis, ferramentas, \textit{streaming} e contabilização de \textit{tokens}, mas implementam-nos de formas estruturalmente diferentes. As diferenças manifestam-se na nomenclatura dos campos, na organização hierárquica do JSON, na localização dos parâmetros de configuração e nos mecanismos de ativação de funcionalidades como \textit{streaming} e raciocínio. Contudo, apesar destas similaridades conceituais, ainda existem campos proprietários em cada API de fornecedor de inferência. Esta análise permitiu perceber que uma abstração que pretenda cobrir estes três fornecedores precisa de ir além de uma simples tradução de nomes, exigindo uma compreensão das diferenças semânticas e estruturais observadas nos pedidos e respostas reais.
% -----------------------------------------------------------------------
\section{O protocolo MCP}
% -----------------------------------------------------------------------

O \textbf{Model Context Protocol (MCP)} é um protocolo aberto desenvolvido pela
Anthropic com o objetivo de normalizar a forma como aplicações baseadas em modelos
de linguagem interagem com ferramentas, fontes de dados e outros sistemas externos
\cite{mcp_protocol}. O protocolo define uma interface comum entre aplicações de IA
e serviços externos, permitindo que diferentes modelos e clientes utilizem as mesmas
capacidades sem necessidade de desenvolver integrações específicas para cada
fornecedor.

Antes do surgimento do \textbf{MCP}, cada plataforma disponibilizava mecanismos
próprios para permitir que os modelos invocassem ferramentas externas. Embora estes
mecanismos oferecessem funcionalidades semelhantes, apresentavam interfaces
incompatíveis entre si, obrigando os programadores a implementar adaptações
específicas para cada API. Esta situação dificultava a reutilização de integrações,
aumentava o esforço de desenvolvimento e criava um forte acoplamento entre as
aplicações e os fornecedores dos modelos.

Para além da uniformização das integrações, o protocolo procura resolver outra
limitação inerente aos modelos de linguagem: o acesso a informação externa. Um
modelo de linguagem possui apenas o conhecimento adquirido durante o processo de
treino e não tem acesso direto a informação dinâmica, como bases de dados, sistemas
empresariais, ficheiros locais ou serviços externos. O \textbf{MCP} permite que
essas capacidades sejam disponibilizadas através de uma interface padronizada,
possibilitando que o modelo obtenha contexto atualizado ou execute operações no
exterior sempre que necessário.

A adoção do \textbf{MCP} introduz assim uma camada de abstração independente dos
fornecedores de modelos de linguagem, uniformizando o acesso a ferramentas, recursos
e outras fontes de contexto externo. Enquanto as APIs de inferência definem a
comunicação entre a aplicação e os modelos, o \textbf{MCP} concentra-se
exclusivamente na interação com sistemas externos, promovendo uma clara separação de
responsabilidades que se revelou determinante na conceção da arquitetura da
\textbf{OmniAI}.

% -----------------------------------------------------------------------
\subsection{Arquitetura do protocolo}
% -----------------------------------------------------------------------

O protocolo segue uma arquitetura cliente-servidor. Nesta arquitetura, os servidores
disponibilizam capacidades que podem ser utilizadas pelos modelos, enquanto os
clientes estabelecem ligações aos servidores, descobrem as capacidades disponíveis e
efetuam a respetiva utilização. Ao contrário das APIs de inferência, responsáveis
pela comunicação entre a aplicação e o modelo de linguagem, o \textbf{MCP} define
exclusivamente a comunicação entre aplicações e fontes externas de contexto. Esta
separação permite desacoplar o acesso a ferramentas da interação com os modelos,
tornando ambas as componentes independentes.

A Figura~\ref{fig:mcp_architecture} ilustra esta separação de responsabilidades.
Do lado da inferência, a aplicação comunica diretamente com o modelo de linguagem
através de uma API de inferência. Do lado do contexto externo, o cliente \textbf{MCP}
estabelece ligação a um ou mais servidores \textbf{MCP}, cada um responsável por
disponibilizar um conjunto de capacidades específicas.

\begin{figure}[h]
\centering
\begin{tikzpicture}[
  node distance = 1.8cm and 2.5cm,
  box/.style     = {rectangle, rounded corners=4pt, minimum width=2.8cm,
                    minimum height=1.0cm, align=center, font=\small\sffamily},
  lbl/.style     = {font=\scriptsize\sffamily\itshape, text=mcpgray},
  arr/.style     = {-stealth, thick},
  darr/.style    = {stealth-stealth, thick},
]

% ---- Application (centre-left) ----
\node[box, fill=mcpvlight, draw=mcpdark, text=mcpdark, thick] (app)
      {Aplicação};

% ---- LLM (top-right) ----
\node[box, fill=mcplight, draw=mcpgray, text=mcpdark,
      above right=1.2cm and 2.5cm of app] (llm)
      {Modelo de\\linguagem};

% ---- MCP Client (bottom, inside app conceptually) ----
\node[box, fill=mcpvlight, draw=mcpgray, text=mcpdark,
      below right=1.2cm and 2.5cm of app] (client)
      {Cliente MCP};

% ---- MCP Servers ----
\node[box, fill=mcplight, draw=mcpdark, text=mcpdark,
      right=3.0cm of client, yshift= 0.9cm] (srv1)
      {Servidor MCP\\(base de dados)};
\node[box, fill=mcplight, draw=mcpdark, text=mcpdark,
      right=3.0cm of client, yshift=-0.9cm] (srv2)
      {Servidor MCP\\(API externa)};

% ---- Arrows ----
\draw[darr, color=mcpgray]
      (app) -- node[above=2pt, lbl, sloped]{API de inferência} (llm);

\draw[darr, color=mcpdark]
      (app) -- node[below=2pt, lbl, sloped]{integração interna} (client);

\draw[darr, color=mcpdark]
      (client) -- node[above=2pt, lbl, sloped]{MCP} (srv1);
\draw[darr, color=mcpdark]
      (client) -- node[below=2pt, lbl, sloped]{MCP} (srv2);

% ---- Divider between the two "worlds" ----
\draw[dashed, color=mcpborder, thick]
      ($(app.east)!0.45!(client.west) + (0,2.4)$)
   -- ($(app.east)!0.45!(client.west) + (0,-2.4)$);

\node[lbl, rotate=90] at ($(app.east)!0.45!(client.west) + (-0.35,0)$)
      {inferência};
\node[lbl, rotate=90] at ($(app.east)!0.45!(client.west) + (0.45,0)$)
      {contexto externo};

\end{tikzpicture}
\caption{Arquitetura do MCP: separação entre a camada de inferência (API do modelo)
         e a camada de contexto externo (servidores MCP).}
\label{fig:mcp_architecture}
\end{figure}

As capacidades disponibilizadas pelos servidores encontram-se organizadas em três
categorias principais, ilustradas na Figura~\ref{fig:mcp_capabilities}:

\begin{itemize}
  \item \textbf{Tools}: operações que podem ser executadas pelo cliente, permitindo
        realizar ações sobre sistemas externos, como executar comandos, consultar
        bases de dados ou interagir com APIs.
  \item \textbf{Resources}: informação disponibilizada pelo servidor para consulta,
        normalmente identificada através de URIs, incluindo documentos, ficheiros de
        configuração ou outros dados estruturados.
  \item \textbf{Prompts}: modelos de interação reutilizáveis que podem ser
        parametrizados pelo cliente para diferentes cenários de utilização.
\end{itemize}

\begin{figure}[h]
\centering
\begin{tikzpicture}[
  every node/.style={font=\small\sffamily},
  cat/.style  = {rectangle, rounded corners=6pt, minimum width=3.2cm,
                 minimum height=1.8cm, align=center, text width=3.0cm},
  icon/.style = {font=\large},
  lbl/.style  = {font=\footnotesize\sffamily, text=mcpgray},
]

% Server box
\node[rectangle, rounded corners=8pt, draw=mcpgray, fill=mcpvlight,
      minimum width=11.6cm, minimum height=3.8cm] (server) {};
\node[lbl, anchor=north west] at (server.north west) {\quad Servidor MCP};

% Three categories inside
\node[cat, fill=white, draw=mcpdark, text=mcpdark, thick,
      left=0.4cm of server.center, xshift=-1.6cm] (tools)
      {\textbf{Tools}\\[4pt]{\scriptsize\color{mcpgray}Executam ações\\em sistemas externos}};

\node[cat, fill=white, draw=mcpdark, text=mcpdark, thick,
      at=(server.center)] (resources)
      {\textbf{Resources}\\[4pt]{\scriptsize\color{mcpgray}Disponibilizam\\informação via URI}};

\node[cat, fill=white, draw=mcpdark, text=mcpdark, thick,
      right=0.4cm of server.center, xshift=1.6cm] (prompts)
      {\textbf{Prompts}\\[4pt]{\scriptsize\color{mcpgray}Modelos de interação\\reutilizáveis}};

\end{tikzpicture}
\caption{As três categorias de capacidades disponibilizadas por um servidor MCP.}
\label{fig:mcp_capabilities}
\end{figure}

O cliente não necessita de conhecer previamente quais destas capacidades existem.
Após estabelecer ligação ao servidor, pode solicitar dinamicamente a lista de
ferramentas, recursos e \textit{prompts} disponibilizados, adaptando automaticamente
o seu comportamento às funcionalidades oferecidas pelo servidor.

% -----------------------------------------------------------------------
\subsection{Descoberta de capacidades}
% -----------------------------------------------------------------------

Uma das principais características do protocolo consiste na descoberta dinâmica de
capacidades (\textit{capability discovery}). Depois de estabelecida uma sessão entre
cliente e servidor, o cliente inicia o processo de inicialização através de mensagens
\textbf{JSON-RPC}. Durante esta fase são negociadas as capacidades suportadas por
ambas as partes e estabelecidos os parâmetros da sessão.

Após a inicialização, o cliente pode solicitar ao servidor a lista de ferramentas,
recursos e \textit{prompts} disponíveis. Esta descoberta é feita através de três
operações definidas pelo protocolo: \texttt{tools/list}, \texttt{resources/list} e
\texttt{prompts/list}. A invocação efetiva de uma ferramenta ou a leitura de um
recurso são depois realizadas pelas operações correspondentes, \texttt{tools/call} e
\texttt{resources/read}, respetivamente. Desta forma, diferentes servidores podem
disponibilizar conjuntos distintos de funcionalidades sem que seja necessário alterar
a lógica do cliente.

Quando o modelo necessita de executar uma determinada ação, o cliente seleciona a
ferramenta apropriada, envia os parâmetros necessários para o servidor e recebe
posteriormente o resultado da operação. A Figura~\ref{fig:mcp_sequence} representa o
fluxo completo de comunicação, desde o estabelecimento da sessão até à execução de
uma ferramenta.

\begin{figure}[h]
\centering
\begin{tikzpicture}[
  every node/.style={font=\small\sffamily},
  life/.style  = {draw=mcpborder, thick, dashed},
  head/.style  = {rectangle, rounded corners=4pt, draw, fill, text=white,
                  minimum width=2.6cm, minimum height=0.7cm, align=center,
                  font=\small\sffamily\bfseries},
  msg/.style   = {-stealth, thick},
  rmsg/.style  = {stealth-, thick, dashed},
  lbl/.style   = {font=\scriptsize\sffamily, fill=white, inner sep=1pt},
  grp/.style   = {rectangle, rounded corners=3pt, draw=mcpborder,
                  fill=mcplight, inner sep=4pt},
]

% ---- Lifeline headers ----
\node[head, fill=mcpdark]   (CH) at (0,0)   {Cliente MCP};
\node[head, fill=mcpgray]   (SH) at (8,0)   {Servidor MCP};

% ---- Lifelines ----
\draw[life] (CH.south) -- ++(0,-9.6);
\draw[life] (SH.south) -- ++(0,-9.6);

% ---- Helper macro positions (y offsets) ----
% 1. initialize  -1.0
% 2. initialized -1.8
% 3. tools/list  -2.8
% 4. list result -3.6
% 5. tools/call  -5.0  (after loop label)
% 6. call result -5.8
% 7. (loop end)  -7.0

% -- Phase label: Inicialização --
\node[lbl, text=mcpdark, anchor=west] at (-2.2,-0.85) {\textit{Inicialização}};
\draw[thick, color=mcpborder] (-2.2,-0.7) -- (-2.2,-2.5);
\draw[thick, color=mcpborder] (-2.2,-0.7) -- (-2.1,-0.7);
\draw[thick, color=mcpborder] (-2.2,-2.5) -- (-2.1,-2.5);

\draw[msg, color=mcpdark]
      (0,-1.0) -- node[lbl, above]{\texttt{initialize}} (8,-1.0);
\draw[rmsg, color=mcpgray]
      (0,-1.8) -- node[lbl, above]{\texttt{initialized} (capacidades negociadas)} (8,-1.8);

% -- Phase label: Descoberta --
\node[lbl, text=mcpdark, anchor=west] at (-2.2,-2.65) {\textit{Descoberta}};
\draw[thick, color=mcpborder] (-2.2,-2.5) -- (-2.2,-4.5);
\draw[thick, color=mcpborder] (-2.2,-4.5) -- (-2.1,-4.5);

\draw[msg, color=mcpdark]
      (0,-2.8) -- node[lbl, above]{\texttt{tools/list}} (8,-2.8);
\draw[rmsg, color=mcpgray]
      (0,-3.6) -- node[lbl, above]{lista de ferramentas disponíveis} (8,-3.6);

% -- Phase label: Execução --
\node[lbl, text=mcpdark, anchor=west] at (-2.2,-4.85) {\textit{Execução}};
\draw[thick, color=mcpborder] (-2.2,-4.7) -- (-2.2,-6.7);
\draw[thick, color=mcpborder] (-2.2,-6.7) -- (-2.1,-6.7);

% loop box frame (fill=none para não tapar as lifelines)
\draw[thick, color=mcpgray, rounded corners=2pt] (-0.8,-4.7) rectangle (8.8,-6.5);
\node[draw=mcpgray, thick, fill=white, anchor=north west, rounded corners=1pt, font=\scriptsize\sffamily\bfseries, inner sep=3pt] at (-0.8,-4.7) {loop};
\node[font=\scriptsize\sffamily, text=mcpgray, anchor=west] at (0.2,-4.95) {[para cada ferramenta invocada]};

\draw[msg, color=mcpdark]
      (0,-5.4) -- node[lbl, above]{\texttt{tools/call} (nome + parâmetros)} (8,-5.4);
\draw[rmsg, color=mcpgray]
      (0,-6.2) -- node[lbl, above]{resultado da operação} (8,-6.2);

\end{tikzpicture}
\caption{Sequência de comunicação MCP: inicialização da sessão, descoberta de
         capacidades e execução de ferramentas via JSON-RPC.}
\label{fig:mcp_sequence}
\end{figure}

Este mecanismo permite que novas funcionalidades sejam adicionadas aos servidores sem
modificar a lógica das aplicações clientes.

% -----------------------------------------------------------------------
\subsection{Mecanismos de transporte}
% -----------------------------------------------------------------------

O protocolo define o formato das mensagens e o ciclo de comunicação entre cliente e
servidor, mas permite que essa comunicação seja realizada através de diferentes
mecanismos de transporte. A versão atual da especificação define dois mecanismos
oficiais: \textbf{STDIO} e \textbf{Streamable HTTP}
\cite{mcp_architecture}, cujas diferenças estão resumidas na
Tabela~\ref{tab:mcp_transport}.

\begin{table}[H]
\centering
\small
\renewcommand{\arraystretch}{1.5}
\begin{tabularx}{\textwidth}{|>{\raggedright\arraybackslash}p{3cm}|Y|Y|}
\hline
\textbf{Característica} & \textbf{STDIO} & \textbf{Streamable HTTP} \\
\hline
\textbf{Âmbito} & Processos locais & Processos distribuídos \\
\hline
\textbf{Transporte} & \texttt{stdin} / \texttt{stdout} & HTTP \\
\hline
\textbf{Streaming} & --- & Suportado \\
\hline
\textbf{Caso de uso típico} & Aplicações desktop (ex.: Claude Desktop) & Servidores remotos e arquiteturas cloud \\
\hline
\end{tabularx}
\caption{Comparação entre os dois mecanismos de transporte oficiais do MCP.}
\label{tab:mcp_transport}
\end{table}

O transporte \textbf{STDIO} destina-se principalmente à comunicação entre processos
executados localmente. Neste modo, o cliente inicia o processo correspondente ao
servidor \textbf{MCP} e toda a comunicação é realizada através dos fluxos de entrada
e saída padrão (\textit{stdin} e \textit{stdout}). Esta abordagem apresenta uma
implementação simples e elimina a necessidade de configurar ligações de rede, sendo
por isso amplamente utilizada por aplicações de ambiente de trabalho como o Claude
Desktop.

O transporte \textbf{Streamable HTTP} destina-se à comunicação entre processos
distribuídos através da rede. Neste caso, as mensagens JSON-RPC são transportadas
sobre HTTP, sendo igualmente possível transmitir respostas de forma incremental
(\textit{streaming}) quando necessário. Este mecanismo substitui a abordagem
anteriormente baseada em \textbf{HTTP} e \textbf{Server-Sent Events (SSE)},
oferecendo uma solução mais simples e flexível para comunicações remotas.

Embora a especificação atual privilegie estes dois mecanismos, diversas
implementações disponibilizam suporte adicional para outros transportes,
nomeadamente \textbf{WebSocket}. Apesar de não constituir um mecanismo oficial do
protocolo, o WebSocket é frequentemente utilizado em aplicações que necessitam de
comunicação bidirecional persistente e de baixa latência.

% -----------------------------------------------------------------------
\subsection{Ferramentas de inspeção}
% -----------------------------------------------------------------------

Durante o desenvolvimento do protocolo foram igualmente disponibilizadas ferramentas
que facilitam a validação de implementações. Entre estas destaca-se o \textbf{MCP
Inspector}, distribuído através do ecossistema Node.js.

Esta ferramenta permite estabelecer ligação a um servidor \textbf{MCP}, visualizar
interativamente as capacidades disponibilizadas, invocar ferramentas, consultar
recursos e analisar as mensagens JSON-RPC trocadas durante toda a sessão. A
utilização do \textbf{MCP Inspector} revelou-se particularmente útil durante o
desenvolvimento do protótipo, permitindo validar o comportamento do servidor e
confirmar a conformidade da implementação com a especificação do protocolo.

% -----------------------------------------------------------------------
\subsection{Protótipo desenvolvido}
% -----------------------------------------------------------------------

Para complementar o estudo da especificação foi desenvolvido um protótipo utilizando
a biblioteca oficial \texttt{io.modelcontextprotocol:kotlin-sdk}
\cite{mcp_kotlin_sdk}. O objetivo deste protótipo consistiu em compreender o ciclo
completo de funcionamento do protocolo e avaliar a maturidade da biblioteca para
futura integração na arquitetura da \textbf{OmniAI}.

O trabalho realizado incluiu a implementação de um servidor \textbf{MCP}, responsável
pelo registo de ferramentas e restantes capacidades disponibilizadas, e de um cliente
\textbf{MCP}, responsável por estabelecer ligação ao servidor, realizar o processo de
inicialização, descobrir dinamicamente as capacidades disponíveis e invocar as
ferramentas registadas.

Durante o desenvolvimento foram explorados os diferentes mecanismos de transporte
suportados pela biblioteca, tendo sido igualmente utilizado o \textbf{MCP Inspector}
para validar o correto funcionamento da comunicação entre cliente e servidor e
verificar a troca de mensagens definida pela especificação.

Este protótipo permitiu compreender detalhadamente o funcionamento do protocolo,
desde o estabelecimento da sessão até à descoberta de capacidades e execução de
ferramentas, fornecendo uma experiência prática complementar ao estudo da
documentação oficial.

\section{Estado da Arte de Gateways de IA e Soluções Semelhantes}

No momento em que o desenvolvimento deste projeto teve início, a adoção de componentes dedicados ao intermédio entre aplicações e Large Language Models (LLMs) era ainda uma prática pouco comum. Existiam escassas arquiteturas formalizadas e os produtos disponíveis no mercado com o propósito de atuar como \textit{gateways} de IA eram raros. A grande parte das aplicações integrava as APIs dos fornecedores de forma direta, o que de certa forma acabava por limitar a escalabilidade e a gestão centralizada.

Com a evolução exponencial do uso de modelos de linguagem em ambientes de produção, surgiu uma necessidade de governança, de controlo de custos, de segurança e de observabilidade. No decorrer do nosso desenvolvimento, o cenário de mercado transformou-se por completo. Atualmente, existem centenas de soluções e arquiteturas adotadas por empresas que integram nativamente uma \textit{gateway} de IA nos seus sistemas. Surgiram várias abordagens que visam resolver os mesmos problemas base, mas com focos distintos:

\begin{itemize}
    \item \textbf{LiteLLM} \cite{litellm}: Uma solução amplamente adotada que atua como um \textit{proxy} central para unificar chamadas a dezenas de fornecedores utilizando o formato da \textbf{OpenAI}. Destaca-se por fornecer funcionalidades prontas de balanceamento de carga e controlo de orçamentos.
    \item \textbf{RouteLLM} \cite{routellm}: Uma \textit{framework} focada no roteamento inteligente de LLMs. A sua principal proposta de valor é a otimização de custo e latência, reencaminhando pedidos mais simples para modelos mais baratos e rápidos, e reservando os modelos de maior capacidade para tarefas mais complexas.
    \item \textbf{Cloudflare AI Gateway} \cite{cloudflare_ai_gateway}: Integrado na rede global da Cloudflare, este serviço posiciona a gestão das chamadas na \textit{edge}, proporcionando benefícios significativos a nível de latência e de \textit{caching} de respostas sem exigir gestão de infraestrutura.
    \item \textbf{Kong AI Gateway} \cite{kong_ai_gateway}: Focado em ambientes empresariais robustos, expande os conceitos tradicionais de gestão de APIs (como \textit{rate limiting} e autenticação) para cobrir o tráfego de inteligência artificial, oferecendo políticas de segurança mais rígidas e suporte para ambientes complexos.
    \item \textbf{Portkey} \cite{portkey_ai}: Uma plataforma com um foco acentuado na observabilidade e operações (\textit{AI Ops}), fornecendo gestão profunda de \textit{prompts}, gestão de \textit{fallbacks} automáticos e análise semântica.
\end{itemize}

Para além destas soluções abertas ou comerciais, a necessidade deste tipo de componente provou ser tão grande que várias das principais empresas de base tecnológica optaram por construir as suas próprias soluções internas. Um caso notável, que demonstra a centralidade de uma \textit{gateway} de IA em arquiteturas de larga escala, foi reportado pela \textbf{Uber} \cite{uber_ai_gateway}. A empresa construiu a sua própria arquitetura interna no interior da sua plataforma de \textit{machine learning} (Michelangelo) como forma de resolver a fragmentação de acessos e garantir um controlo de segurança unificado. Isto tornou-se crucial num ecossistema em que múltiplos agentes de IA distintos necessitam de interagir em harmonia, requerendo uma forte gestão centralizada de identidades e permissões.

A constatação de que o mercado validou a necessidade de um componente com as mesmas responsabilidades que a nossa \textbf{OmniAI} reforça a pertinência do problema escolhido. A adoção massiva destas tecnologias prova que o caminho centralizado é a direção tomada pela indústria. O próximo capítulo irá detalhar como a arquitetura do nosso projeto responde de forma inovadora a estes mesmos desafios.